\section{Sampling starvation and robust optimisation}
\label{sec:sampling}

\textbf{Sampling starvation.} The surface term is a Monte Carlo estimate,
$\mathcal{L}_{\mathrm{ch}}\approx\frac1N\sum_i d(p_i,T)^2$ with
$p_i\sim p(x)$. If $p$ is proportional to surface area, the expected loss
$\mathbb{E}_{x\sim p}[\mathrm{err}(x)]$ weights the body and head far more
than distal legs or antennae, so a tiny structure can be almost invisible to
the optimiser: a second route to the under-representation of \F{F5}, this
time in the sampling scheme rather than the mesh. Reserving a fixed quota of
the global budget for distal points improved distal legs but starved
antennae and head, which shared that budget; restricting the quota to within
the leg budget avoided the trade-off, yet did not clearly improve on simply
splitting distal sampling. The conclusion is narrow: starvation is a real
confound, but controlling it did not create anatomical correspondence.

\textbf{Robust optimisation.} A related question is how the objective
should respond to points that are simply wrong (occluded regions,
scanning artefacts, disconnected debris) rather than merely
under-sampled. \term{Graduated non-convexity} (GNC)~\cite{gnc} is a robust-optimisation
strategy that treats every point correspondence as suspect initially and
progressively hardens its confidence in inliers as optimisation proceeds,
rather than trusting every nearest-point match from the first iteration.
Applied globally across the whole body, it produced undesirable behaviour
specifically around the antennae; restricted to the leg point set alone
(\emph{leg-only, topology-free}), it improved robustness where applied,
under synthetic corruption benchmarks with 30\% and 60\% of points dropped:
accuracy rose from $0.602$ to $0.797$ (Drop30) and from $0.411$ to $0.607$
(Drop60). Localising the robust treatment to the anatomical subsystem where
corrupted points actually occur outperformed applying it uniformly; the
result is retained as a promising candidate for leg registration
specifically, not adopted as a general replacement for point-matching across
the whole body.
